\documentclass[11pt,a4paper]{article}

% ---------- Datos del documento ----------
\newcommand{\titulo}{Caso 9. Adolescente con problemas de peso: análisis DAFO y líneas de actuación}
\newcommand{\tituloCorto}{Caso 9}
\newcommand{\asignatura}{Organización y Gestión del Centro}
\newcommand{\autor}{Javier Gómez López}
\newcommand{\enunciado}{%
A partir del caso \emph{Adolescente con problemas de peso}:
\begin{enumerate}
\item Analizar el caso y elaborar un análisis DAFO.
\item Proponer líneas de actuación que aporten soluciones para mejorar la situación de la adolescente en el centro educativo.
\item Ordenar las actuaciones por su grado de urgencia e importancia.
\item Clasificar las actuaciones según los órganos, departamentos o personal del centro responsables de ellas.
\end{enumerate}
}
\newcommand{\master}{MÁSTER UNIVERSITARIO EN FORMACIÓN DEL PROFESORADO-ESPECIALIDAD EN MATEMÁTICAS -- Universidad Rey Juan Carlos}

% ---------- Paquetes ----------
\usepackage[utf8]{inputenc}
\usepackage[T1]{fontenc}
\usepackage[spanish,es-noshorthands]{babel}
\usepackage{lmodern}
\usepackage[margin=2.5cm,headheight=14pt]{geometry}
\usepackage{xcolor}
\usepackage{titlesec}
\usepackage{fancyhdr}
\usepackage{enumitem}
\usepackage{parskip}
\usepackage{etoolbox}
\usepackage{tabularx}
\usepackage{booktabs}
\usepackage{longtable}
\usepackage{array}
\usepackage[most]{tcolorbox}
\usepackage[hidelinks]{hyperref}

% ---------- Estilo ----------
\definecolor{acento}{RGB}{176,0,32}   % rojo URJC
\definecolor{gris}{RGB}{90,90,90}
\definecolor{verde}{RGB}{46,125,50}
\definecolor{azul}{RGB}{21,101,192}
\definecolor{naranja}{RGB}{230,120,0}

\titleformat{\section}{\Large\bfseries\color{acento}}{\thesection.}{0.6em}{}
\titleformat{\subsection}{\large\bfseries\color{gris}}{\thesubsection.}{0.6em}{}

\pagestyle{fancy}
\fancyhf{}
\fancyhead[L]{\small\color{gris}\asignatura}
\fancyhead[R]{\small\color{gris}\tituloCorto}
\fancyfoot[C]{\small\thepage}
\renewcommand{\headrule}{{\color{acento}\hrule height 0.6pt}}

\newcommand{\dafobox}[3]{%
  \begin{tcolorbox}[enhanced, colback=#1!5, colframe=#1, boxrule=0.8pt, arc=1pt,
      title=#2, fonttitle=\bfseries, height=9.6cm, valign=top, fontupper=\small]
    #3
  \end{tcolorbox}}
\newcommand{\prio}[1]{\textbf{\color{acento}#1}}

\begin{document}

% ---------- Cabecera ----------
\thispagestyle{empty}
\begin{center}
  {\small\color{gris}\MakeUppercase{\master}}\\[1.2em]
  {\color{acento}\rule{\linewidth}{1.2pt}}\\[0.8em]
  {\Huge\bfseries \titulo}\\[0.6em]
  {\Large\color{gris} \asignatura}\\[0.6em]
  {\color{acento}\rule{\linewidth}{1.2pt}}\\[1em]
  {\large \autor}\\[0.3em]
  {\small\color{gris}\today}
\end{center}
\vspace{1.5em}

% ---------- Enunciado ----------
\ifdefempty{\enunciado}{}{%
  \begin{tcolorbox}[enhanced, breakable, colback=acento!4, colframe=acento,
      boxrule=0pt, leftrule=3pt, arc=0pt, title=Enunciado de la tarea,
      coltitle=acento, colbacktitle=acento!4, fonttitle=\bfseries]
    \enunciado
  \end{tcolorbox}
  \vspace{1em}
}

% ---------- Contenido ----------
\section{Análisis del caso}

\subsection{Datos relevantes}
\begin{description}[leftmargin=0pt, style=nextline]
  \item[Situación sanitaria] Mujer de 20 años (adolescencia tardía) que acude
  al centro de salud con su madre por delgadez, falta de apetito y comidas a
  deshoras. Peso de 48\,kg, talla de 160\,cm e IMC de 18,75, en el límite
  inferior de la normalidad (18,5). Presenta palidez cutánea y el resto de la
  exploración es normal. No toma medicación. Además, la madre solicita una
  serología de VIH.
  \item[Conductas de riesgo] Se marchó de casa tres días con una amiga
  consumidora de drogas y alcohol, después de robar a su padre 60.000\,€ que el
  grupo gastó en ese tiempo (entre otras cosas, en un coche). Según la madre,
  hay riesgo de accidentes, de consumo de drogas y alcohol, de embarazo y de
  delincuencia. Solo lo contó después de mucha insistencia.
  \item[Situación escolar] Solo ha cursado hasta 2.º de ESO y no tiene el título
  de Graduado en ESO. Hace dos años que dejó el instituto porque, en sus
  palabras, \emph{<<no le interesaba nada de lo que se podía aprender allí>>}.
  Es un caso de abandono escolar temprano. Como ya ha superado la edad de la
  enseñanza obligatoria, hoy está fuera del sistema educativo.
  \item[Contexto familiar (genograma)] El padre tiene 43 años, estudios de EGB y
  trabaja como camionero. El genograma lo señala como alcohólico, con esguince
  cervical y lumbalgia, y la relación con su hija aparece como distante. La
  madre tiene 41 años, estudios universitarios y es profesora; el genograma
  recoge la anotación <<DP>> (previsiblemente, depresión). Las relaciones de la
  madre con el padre y con la hija se representan como conflictivas. No
  constan hermanos.
\end{description}

\subsection{Interpretación desde el centro educativo}
El problema que motiva la consulta es el bajo peso, pero desde el punto de vista
educativo es un síntoma más de una situación de \textbf{riesgo psicosocial
múltiple}. Coinciden aquí la desvinculación escolar, los hábitos de salud
inadecuados, el consumo de sustancias en el grupo de iguales, las conductas
antisociales y un entorno familiar con alcoholismo, posible patología depresiva
y relaciones conflictivas.

Que tenga \textbf{20 años} cambia el planteamiento de la intervención:
\begin{itemize}
  \item \textbf{El centro de referencia ya no es el instituto ordinario.} Es un
  \textbf{centro de educación de personas adultas} (CEPA), o un IES que ofrezca
  educación de adultos, al que puede incorporarse porque ya tiene 18 años
  (art.~67.1 LOMLOE). Ahí puede cursar la Educación Secundaria para Personas
  Adultas (ESPA), presentarse a las pruebas libres para obtener el título de
  la ESO (art.~68.2) o acceder a un ciclo formativo de grado básico para
  adultos (art.~68.3). La LOMLOE incluye además entre los objetivos de la
  educación de adultos <<promover la salud y los hábitos saludables de
  alimentación>> (art.~66.3.i).
  \item \textbf{No hay absentismo ni protección de menores.} Su vuelta a la
  formación es \textbf{voluntaria}: no se puede obligar, hay que motivarla.
  Por eso la acogida y el vínculo inicial son decisivos.
  \item \textbf{Ella decide sobre sí misma.} Es mayor de edad: la información
  sanitaria y la implicación de la familia en el centro requieren \textbf{su
  consentimiento}. La madre deja de ser la interlocutora formal, aunque sigue
  siendo un apoyo muy valioso.
  \item \textbf{Las consecuencias de sus actos son de adulta.} Un robo o un
  delito ya no se tratan por la vía de menores, lo que hace más urgente
  intervenir antes de que se cronifique la situación.
\end{itemize}

Por tanto, el análisis se sitúa en el \textbf{centro de adultos donde la joven
se matricula}. En este caso, la derivación puede venir del propio centro de
salud o de su madre.

\section{Análisis DAFO}
Los factores internos son los de la propia joven y el centro de adultos. Los
externos, los de la familia, el grupo de iguales y los recursos del entorno.

\noindent
\begin{minipage}[t]{0.49\linewidth}
\dafobox{acento}{Debilidades (internas)}{%
\begin{itemize}[leftmargin=1em, itemsep=1pt]
  \item Desmotivación hacia lo escolar y dos años desconectada de cualquier
  formación.
  \item Sin título de la ESO y con un nivel formativo bajo (2.º de ESO).
  \item Bajo peso, inapetencia, comidas desordenadas y palidez: posible riesgo de
  trastorno de la conducta alimentaria o de malnutrición asociada al consumo.
  \item Conductas de riesgo: fuga, robo, posible consumo de alcohol y drogas, y
  relaciones sexuales de riesgo (VIH, embarazo).
  \item Baja percepción del riesgo y comunicación difícil con los adultos.
  \item Centro: alto riesgo de que la deje pronto. El alumnado adulto se da de
  baja con facilidad si no se siente acogido.
\end{itemize}}
\end{minipage}\hfill
\begin{minipage}[t]{0.49\linewidth}
\dafobox{naranja}{Amenazas (externas)}{%
\begin{itemize}[leftmargin=1em, itemsep=1pt]
  \item Al no ser obligatoria la escolarización, ningún mecanismo la hace
  volver si ella no quiere.
  \item Padre con alcoholismo y relación distante con la hija; relaciones
  conflictivas en la familia; posible depresión materna.
  \item Grupo de iguales consumidor de drogas y alcohol, con acceso a dinero.
  \item Consecuencias penales como adulta, riesgo de accidentes, de embarazo no
  deseado y de infecciones de transmisión sexual.
  \item Con 20 años y sin título de la ESO, riesgo de exclusión laboral y de
  empleo precario.
\end{itemize}}
\end{minipage}

\noindent
\begin{minipage}[t]{0.49\linewidth}
\dafobox{verde}{Fortalezas (internas)}{%
\begin{itemize}[leftmargin=1em, itemsep=1pt]
  \item Buen estado físico general: consciente, orientada, colaboradora,
  exploración normal y sin medicación.
  \item Acude a consulta y acaba contando lo ocurrido: hay margen para el
  vínculo y la comunicación.
  \item Mayoría de edad: puede decidir y elegir su itinerario (ESPA, pruebas
  libres, FP), que puede tener una orientación profesional.
  \item Su desinterés apunta a una oferta que no conecta con ella, no
  necesariamente a una falta de capacidad.
  \item Centro de adultos: horarios flexibles, modalidades presencial y a
  distancia (art.~70 bis), grupos reducidos y enseñanza organizada por ámbitos.
\end{itemize}}
\end{minipage}\hfill
\begin{minipage}[t]{0.49\linewidth}
\dafobox{azul}{Oportunidades (externas)}{%
\begin{itemize}[leftmargin=1em, itemsep=1pt]
  \item Madre muy implicada y preocupada; además es profesora y conoce el
  sistema educativo.
  \item Atención primaria ya implicada: puede derivarla al centro de adultos
  y coordinarse con él.
  \item Recursos externos: servicios sociales, centros de atención a las
  adicciones, salud mental y servicios públicos de empleo.
  \item Marco LOMLOE: acceso a la educación de adultos (art.~67.1), salud y
  alimentación (66.3.i), orientación profesional junto con empleo (67.10),
  pruebas libres de la ESO (68.2), grado básico para adultos (68.3) y prueba de
  acceso a FP de grado medio (art.~41).
\end{itemize}}
\end{minipage}

\subsection*{Estrategias que se derivan del DAFO}
\begin{itemize}
  \item \textbf{Fortalezas + Oportunidades (ofensivas).} Aprovechar su buena
  disposición en consulta y la derivación desde el centro de salud para que se
  matricule en un itinerario que ella misma elija.
  \item \textbf{Fortalezas + Amenazas (defensivas).} Usar la tutoría y la
  flexibilidad del centro de adultos para ofrecerle un adulto de referencia y
  un grupo alternativo al grupo consumidor.
  \item \textbf{Debilidades + Oportunidades (reorientación).} Compensar la
  desmotivación con una formación práctica orientada al empleo y aprovechar los
  recursos sanitarios para el problema de peso y de consumo.
  \item \textbf{Debilidades + Amenazas (supervivencia).} Hacer una acogida
  inmediata y un seguimiento estrecho para evitar que lo deje, y derivarla con
  su consentimiento a servicios sanitarios y sociales.
\end{itemize}

\section{Líneas de actuación ordenadas por urgencia e importancia}
Las actuaciones se clasifican con la matriz \emph{urgente--importante} en tres
niveles de prioridad:
\begin{itemize}[itemsep=0pt]
  \item \prio{P1. Inmediata}: urgente e importante, entre la matrícula y las
  primeras semanas.
  \item \prio{P2. Corto plazo}: importante y con plazo, durante el primer
  trimestre y después de forma continuada.
  \item \prio{P3. Medio plazo}: importante pero no urgente, durante el curso y
  en los documentos del centro.
\end{itemize}

{\small
\renewcommand{\arraystretch}{1.25}
\begin{longtable}{>{\bfseries}p{0.06\linewidth} p{0.47\linewidth} >{\centering\arraybackslash}p{0.09\linewidth} >{\raggedright\arraybackslash}p{0.26\linewidth}}
\toprule
Nº & Actuación & Prioridad & Responsable principal \\
\midrule
\endhead
\multicolumn{4}{l}{\prio{P1 --- Inmediata (urgente e importante)}}\\
A1 & \textbf{Acoger e informar.} Atender la derivación del centro de salud o de
la madre, hacer una entrevista de acogida individual, explicarle las opciones
(ESPA presencial o a distancia, pruebas libres de la ESO, grado básico para
adultos, prueba de acceso a grado medio) y facilitarle la matrícula en cualquier
momento del curso. & P1 & Jefatura de estudios; Orientador/a; Secretaría \\
A2 & \textbf{Hacer una valoración inicial y una entrevista de orientación.}
Valorar su nivel de competencia para situarla en el nivel y los módulos que le
correspondan, y hacer una entrevista confidencial sobre sus intereses, su
estado emocional, la autoestima y la relación con la comida. & P1 &
Orientador/a; Equipo docente (ámbitos) \\
A3 & \textbf{Coordinarse con los servicios sanitarios y sociales, con su
consentimiento.} Coordinarse con atención primaria por el bajo peso y la salud
sexual, ofrecerle la derivación a recursos de adicciones, salud mental o
servicios sociales y, si aparecen indicios de violencia o explotación, aplicar
el protocolo del centro. & P1 & Dirección; Orientador/a; Coordinador/a de
bienestar y protección (art.~124.5) \\
A4 & \textbf{Asignar un tutor o tutora de referencia y firmar un compromiso.}
Un adulto de referencia y un compromiso personal firmado por ella, con
objetivos a corto plazo alcanzables. & P1 & Jefatura de estudios; Tutor/a \\
\midrule
\multicolumn{4}{l}{\prio{P2 --- Corto plazo (importante, con plazo)}}\\
A5 & \textbf{Diseñar un itinerario formativo personalizado.} Combinar la ESPA
(o la preparación de las pruebas libres) con un objetivo profesional concreto:
la prueba de acceso a FP de grado medio, un ciclo de grado básico o la
orientación al empleo junto con los servicios públicos de empleo (art.~67.10).
& P2 & Orientador/a; Tutor/a; Jefatura de estudios \\
A6 & \textbf{Adaptar la metodología a una persona adulta.} Aprendizaje
funcional y ligado a la vida real, proyectos relacionados con sus intereses
y una evaluación ajustada a su punto de partida. Si lo necesita, compaginarlo
con la modalidad a distancia (art.~70 bis). & P2 & Equipo docente;
Departamentos de ámbito \\
A7 & \textbf{Hacer un seguimiento para que no lo deje.} Registro de
asistencia, contacto inmediato si falta varios días seguidos, tutorías
individuales periódicas y reuniones del equipo docente para ajustar la
respuesta. & P2 & Tutor/a; Equipo docente; Jefatura de estudios \\
A8 & \textbf{Implicar a la familia, solo si ella lo autoriza.} Con su
autorización, entrevistas con la madre para coordinar el apoyo en casa,
mediación en el conflicto madre--hija y orientación de la familia hacia
recursos de terapia familiar o de alcoholismo. & P2 & Tutor/a; Orientador/a \\
\midrule
\multicolumn{4}{l}{\prio{P3 --- Medio plazo (importante, no urgente)}}\\
A9 & \textbf{Desarrollar educación para la salud.} Talleres de alimentación
saludable e imagen corporal, prevención de adicciones y educación
afectivo-sexual (VIH, embarazo), integrados en los ámbitos y en la tutoría,
en línea con el art.~66.3.i. & P3 & Orientador/a; Tutores; Departamento del
ámbito científico-tecnológico; CCP \\
A10 & \textbf{Crear vínculos positivos con el centro.} Actividades culturales,
deportivas y de participación, enseñanzas para el desarrollo personal y
aprendizaje cooperativo, para ofrecerle un grupo alternativo. & P3 & Jefatura
de estudios; Responsable de actividades complementarias \\
A11 & \textbf{Revisar los documentos del centro y formar al profesorado.}
Incorporar al proyecto educativo, a la PGA y al plan de acción tutorial
protocolos de acogida, de prevención del abandono y de detección de trastornos
de la conducta alimentaria y consumo en jóvenes adultos, y ofrecer formación
al claustro sobre estas señales de alerta. & P3 & Equipo directivo; Claustro;
CCP; Consejo Escolar \\
A12 & \textbf{Establecer convenios con entidades externas.} Colaboración
estable con el centro de salud, los servicios sociales, los recursos de
adicciones, los servicios de empleo y ONG (art.~127.j) para derivar en ambos
sentidos. & P3 & Consejo Escolar (fija las directrices); Dirección (las
ejecuta) \\
\bottomrule
\end{longtable}
}

\textbf{Criterio de ordenación.} Primero se asegura que la joven
\emph{llegue y se quede} en el centro y que sus necesidades de salud estén
atendidas (A1--A4). En la educación de adultos la primera semana decide si
continúa, y sin vínculo no hay intervención posible. Después se construye un
itinerario con sentido para ella y se evita que lo deje (A5--A8). Por último,
se actúa sobre la organización para responder mejor a otros casos (A9--A12).
Al ser mayor de edad, rigen siempre su \textbf{autonomía y su consentimiento}:
la información sanitaria y familiar solo se comparte si ella lo autoriza y con
quien la necesita para intervenir.

\section{Clasificación por órganos y personal responsables}
Los órganos de gobierno (Consejo Escolar, Claustro y equipo directivo) son los
que recoge la LOMLOE (arts.~119 y 126--132). Los órganos de coordinación docente
(orientación, departamentos, CCP, tutoría y equipo docente) son los que la ley
menciona en su art.~119. En los centros de adultos los desarrolla la normativa
autonómica de forma análoga al Reglamento Orgánico de los IES\footnote{Real
Decreto 83/1996, de 26 de enero, por el que se aprueba el Reglamento Orgánico de
los Institutos de Educación Secundaria.}, con departamentos organizados por
ámbitos.

{\small
\renewcommand{\arraystretch}{1.3}
\begin{longtable}{>{\raggedright\arraybackslash}p{0.27\linewidth} p{0.53\linewidth} >{\centering\arraybackslash}p{0.12\linewidth}}
\toprule
Órgano / personal & Funciones en este caso & Actuaciones \\
\midrule
\endhead
\multicolumn{3}{l}{\textbf{\color{acento}Órganos colegiados de gobierno}}\\
Consejo Escolar & Propone medidas de estilos de vida saludable y de
convivencia (art.~127.g), fija las directrices de colaboración con entidades
(127.j), aprueba el proyecto educativo y evalúa la PGA (127.a--b). & A11, A12 \\
Claustro del profesorado & Aprueba los aspectos educativos de los planes
(acción tutorial, acogida) y propone la formación del profesorado. & A11 \\
\midrule
\multicolumn{3}{l}{\textbf{\color{acento}Órganos unipersonales (equipo directivo, art.~131)}}\\
Dirección & Garantiza el cumplimiento de la normativa (art.~132.d), impulsa la
colaboración con las familias y con otras instituciones (132.g) y coordina la
derivación a servicios externos. & A3, A11, A12 \\
Jefatura de estudios & Organiza la acogida, la asignación de grupo y tutor, el
seguimiento de la asistencia y las actividades. & A1, A4, A5, A7, A10 \\
Secretaría & Gestiona la matrícula (también fuera de plazo, si la normativa lo
permite) y las inscripciones en pruebas libres, y custodia el expediente de
forma confidencial. & A1 \\
Coordinador/a de bienestar y protección & Aplica el protocolo si aparecen
indicios de violencia o explotación (art.~124.5). & A3 \\
\midrule
\multicolumn{3}{l}{\textbf{\color{acento}Órganos de coordinación docente}}\\
Orientación (departamento u orientador/a del centro) & Hace la entrevista de
acogida, la orientación académico-profesional, la coordinación con los
servicios sanitarios, sociales y de empleo y la intervención familiar
consentida. & A1, A2, A3, A5, A8, A9 \\
Comisión de coordinación pedagógica (CCP) & Integra la educación para la salud
en las programaciones y revisa los documentos institucionales. & A9, A11 \\
Departamentos de ámbito (comunicación, social, científico-tecnológico) & Hacen
la valoración inicial y la adaptación metodológica, y desarrollan contenidos
de salud y alimentación en el ámbito científico-tecnológico. & A2, A6, A9 \\
Responsable de actividades complementarias & Organiza actividades que
favorecen la integración y un ocio saludable. & A10 \\
Tutor/a & Es la persona de referencia: hace el compromiso, las tutorías
individuales, el seguimiento de la asistencia y, si ella lo autoriza, el
contacto con la familia. & A4, A5, A7, A8, A9 \\
Equipo docente & Aplica de forma coherente la adaptación metodológica y
participa en el seguimiento. & A2, A6, A7 \\
\bottomrule
\end{longtable}
}

\section{Conclusión}
El bajo peso de esta joven de 20 años es la parte visible de un problema más
amplio: abandono escolar sin título, conductas de riesgo y un entorno familiar
frágil. Como ya es adulta, el sistema educativo no puede obligarla a volver;
solo puede \textbf{atraerla y retenerla}. Por eso el centro de adultos tiene
que ofrecerle una acogida cálida y un itinerario con sentido profesional, y
actuar de forma coordinada: la jefatura de estudios y la tutoría acogen y hacen
el seguimiento, la orientación evalúa, orienta y conecta con los servicios
sanitarios, sociales y de empleo, el equipo docente adapta la enseñanza a una
persona adulta, y el equipo directivo, el Claustro y el Consejo Escolar
convierten la experiencia en protocolos estables. Todo ello respetando su
autonomía como mayor de edad. La clave es que, mientras siga vinculada a una
formación, sigue habiendo margen para intervenir.

\end{document}
